\chapter{Analiza wymagań i projekt rozwiązania}
Analiza ta pozwala na precyzyjne określenie zadań systemu oraz zaplanowanie jego struktury tak, aby wszystkie komponenty współpracowały bezawaryjnie. Dzięki niej możliwe jest wybranie odpowiednich technologii, które najlepiej sprawdzą się w zamierzonym celu aplikacji.

\section{Wymagania funkcjonalne}
Jak pokazano w tabeli \ref{tab:funkcjonalne}, aplikacja musi spełniać określone funkcje. Każda z nich została przypisana do odpowiedniego modułu, co ułatwia organizację pracy i pozwala na równoległe rozwijanie poszczególnych części projektu.
\begin{table}[ht!]
\centering
\caption{Wymagania funkcjonalne}
\label{tab:funkcjonalne}
\begin{tabular}{lll}
\toprule
ID & Opis & Priorytet \\  
\midrule
FR-001      & Wyświetlanie dzinnika zapisów z datą, godziną i wartościami.    & Wysoki \\
FR-002      & Przygotowanie dokumentacji w oparcu LaTeX. Wymagania edytorskie.    & Wysoki \\
FR-003      & Przygotowanie diagramów UML ilustrujących strukturę i zachowanie aplikacji.    & Średni \\
FR-004      & Utworzenie pliku README.md z opisem projektu i instrukcją instalacji.    & Wysoki \\
\bottomrule 
\end{tabular}
\sourcetab{opracowanie własne}
\end{table}

\section{Wymagania niefunkcjonalne}
Jak pokazano w tabeli \ref{tab:niefunkcjonalne}, aplikacja musi spełniać określone standardy wydajności, bezpieczeństwa i kompatybilności.
\begin{table}[ht!]
\centering
\caption{Wymagania niefunkcjonalne}
\label{tab:niefunkcjonalne}
\begin{tabular}{lll}
\toprule
ID & Opis & Priorytet \\  
\midrule
NFR-001      & Czas odpowiedzi UI mniejszy niż 200ms    & Wysoki \\
NFR-002      & Zużycie baterii mniejszse niż 10\% na godzinę    & Wysoki \\
NFR-003      & Kompatybilność: Android 8.0+	    & Wysoki \\
NFR-004      & Dostępność zgodna z WCAG 2.1 na poziomie AA    & Średni \\
\bottomrule 
\end{tabular}
\sourcetab{opracowanie własne}
\end{table}

\section{Architektura rozwiązania}
Projekt został oparty na sprawdzonej architekturze, która oddziela warstwę danych od warstwy prezentacji. Sercem komunikacji jest biblioteka Retrofit, która w połączeniu z konwerterem Gson pozwala na błyskawiczne przetwarzanie informacji przesyłanych w formacie JSON. Całość oprogramowania została napisana w języku Java.